% Типовой расчёт № 1 по дискретной математике
% Лаборатория математики ФН1
% Компиляция: pdflatex tr-01-usloviya.tex

\documentclass[12pt, a4paper]{article}

\usepackage[T2A]{fontenc}
\usepackage[utf8]{inputenc}
\usepackage[russian]{babel}
\usepackage[left=25mm, right=20mm, top=20mm, bottom=20mm]{geometry}
\usepackage{amsmath, amssymb}
\usepackage{enumitem}

\pagestyle{plain}
\setlength{\parindent}{1.25cm}

\begin{document}

\begin{center}
  {\large\bfseries Типовой расчёт № 1}

  \vspace{0.3em}
  {\large Множества, отношения, графы}

  \vspace{0.5em}
  {\small Дискретная математика · 1 курс, 2 семестр\\
  Лаборатория математики ФН1}
\end{center}

\vspace{1em}

\noindent
Номер варианта совпадает с порядковым номером студента в журнале группы.
Числовые данные вариантов — в файле \texttt{tr-01-varianty.csv}.

\section*{Задача 1. Операции над множествами}

Даны множества
\[
A = \{x \in \mathbb{Z} : |x| \le a\},\quad
B = \{x \in \mathbb{Z} : 0 < x < b\},\quad
C = \{x \in \mathbb{Z} : x = 3k,\ k \in \mathbb{Z}\}.
\]
Найти $(A \cup B) \setminus C$, $A \cap B \cap C$ и симметрическую разность
$A \bigtriangleup B$. Изобразить результат на диаграмме Эйлера~--- Венна.

\section*{Задача 2. Доказательство тождества}

Доказать тождество двумя способами~--- через определение принадлежности
элемента и с помощью диаграмм:
\[
(A \setminus B) \cup (B \setminus A) = (A \cup B) \setminus (A \cap B).
\]

\section*{Задача 3. Бинарные отношения}

На множестве $M = \{1, 2, \dots, n\}$ задано отношение
\[
R = \{(x, y) : x - y \text{ кратно } m\}.
\]
Проверить рефлексивность, симметричность, антисимметричность и транзитивность.
Является ли $R$ отношением эквивалентности? Если да~--- выписать классы
эквивалентности.

\section*{Задача 4. Комбинаторика}

Сколькими способами можно составить команду из $k$ человек, если в группе
$n$ студентов, причём двое из них не могут работать вместе?
Ответ обосновать, приведя схему рассуждения.

\section*{Задача 5. Графы}

Граф задан матрицей смежности (см. файл вариантов). Требуется:
\begin{enumerate}[label=\arabic*)]
  \item изобразить граф;
  \item определить степени всех вершин, проверить лемму о рукопожатиях;
  \item найти все компоненты связности;
  \item выяснить, является ли граф эйлеровым и гамильтоновым;
  \item построить остовное дерево методом поиска в ширину.
\end{enumerate}

\section*{Задача 6. Булевы функции}

Для функции, заданной вектором значений варианта, требуется:
\begin{enumerate}[label=\arabic*)]
  \item построить таблицу истинности;
  \item записать СДНФ и СКНФ;
  \item минимизировать функцию с помощью карты Карно;
  \item проверить принадлежность функции классам
        $T_0$, $T_1$, $S$, $M$, $L$.
\end{enumerate}

\vspace{1.2em}

\section*{Требования к оформлению}

\begin{itemize}[nosep]
  \item все диаграммы и графы вычерчиваются аккуратно, вершины подписываются;
  \item в задаче~4 обязательно словесное обоснование комбинаторной схемы;
  \item в задаче~6 карта Карно приводится полностью, с выделенными областями склеивания.
\end{itemize}

\section*{Критерии оценивания}

\begin{center}
\begin{tabular}{|c|l|}
  \hline
  Баллы & За что начисляются \\
  \hline
  2 & задача решена верно, обоснование полное \\
  1 & верный метод, допущена вычислительная ошибка \\
  0 & метод выбран неверно либо решение отсутствует \\
  \hline
\end{tabular}
\end{center}

\vspace{0.5em}
\noindent
Расчёт зачтён при \textbf{8 баллах из 12} и успешной защите.

\vspace{0.7em}
\noindent\textbf{Сроки.} Выдача~--- 4-я учебная неделя,
сдача~--- \textbf{до конца 10-й недели}.

\end{document}
